\chapter{Степенные ряды. Радиус сходимости. Бесконечная дифференцируемость суммы степенного ряда. Ряд Тейлора.}
\section{Степенные ряды. Радиус сходимости}
\subsubsection{Радиус и круг сходимости степенного ряда. Радиус сходимости}
Рассматриваем функциональные ряды вида
\begin{equation}\label{yaa11p3e1}
\sum\limits_{n=0}^{\infty} a_n(z-z_0)^n,
\end{equation}
где $z_0$ и $a_0,a_1,...,a_n,...$ "--- заданные комплексные числа, а $z$ "--- переменная, принимающая комплексные значения, т.е $z\in\bbC$, где $\bbC$ "--- множество комплексных чисел. Такие ряды называются \textit{степенными рядами}. Числа $a_0,a_1,...,a_n,...$ называются \textit{коэффициентами степенного ряда} \eqref{yaa11p3e1}.

Отметим, что у степенного ряда счет членов ведется не с единицы, а с нуля: первый член называется нулевым, второй "--- первым и т.д.

\begin{thm}[Первая теоремиа Абеля] \label{yaa11p3thm1}
Если степенной ряд \eqref{yaa11p3e1} сходится в точке $z_1\ne z_0$, то он сходится абсолютно в любой точке $z$ из круга $|z-z_0|<r_1$, где $r_1=|z_1-z_0|$.
\end{thm}

\begin{cons}
Если ряд \eqref{yaa11p3e1} расходится в точке $z_2$, то он расходится в любой точке $z$, для которой $|z-z_0|>r_2$, где $r_2=|z_2-z_0|$.
\end{cons}

\begin{defn}
Число $R>0$, обладающее свойством: ряд \eqref{yaa11p3e1} сходится, если $|z-z_0|<R$, и расходится, если $|z-z_0|>R$, называется \textit{радиусом сходимости}, а круг $|z-z_0|<R$ "--- \textit{кругом сходимости} степенного ряда~\eqref{yaa11p3e1}.

Если ряд~\eqref{yaa11p3e1} сходится только в точке $z_0$, то, по определению $R=0$. Если ряд \eqref{yaa11p3e1} сходится при любом $z\in\bbC$, то $R=+\infty$.
\end{defn}

\begin{thm}
У любого степенного ряда~\eqref{yaa11p3e1} существует радиус сходимости $R$, причем
\begin{equation}\label{yaa11p3e2}
R=\frac{1}{\overline{\lim\limits_{n\to\infty}}\sqrt[n]{|a_n|}}.
\end{equation} 
\end{thm}

\begin{proof}
Пусть $\overline{\lim\limits_{n\to\infty}}\sqrt[n]{|a_n|}=q$, и пусть сначала $0<q<+\infty$. Тогда
$$
\overline{\lim\limits_{n\to\infty}}\sqrt[n]{|a_n(z-z_0)^n|}=q\cdot |z-z_0|.
$$
Отсюда по признаку Коши для числового ряда следует, что ряд \eqref{yaa11p3e1} сходится, если $q|z-z_0|<1$, и расходится, если $q|z-z_0|>1$.

Следовательно, $R=1/q$, т.е. справедлива формула~\eqref{yaa11p3e2}.

Если $q=0$, то ряд~\eqref{yaa11p3e1} сходится при любом $z$, и поэтому $R=+\infty$. Если же $q=+\infty$, то ряд \eqref{yaa11p3e1} расходится при любом $z\ne z_0$, и поэтому $R=0$. Можно считать, что в этих случаях тоже справедлива формула~\eqref{yaa11p3e2}. 
\end{proof}

Формула~\eqref{yaa11p3e2} называется \textit{формулой Коши"--~Адамара}\rindex{формула!Коши"---Адамара} для радиуса сходимости степенного ряда~\eqref{yaa11p3e1}.

\begin{thm} \label{yaa11p3thm3}
Степенные ряды
\begin{equation}\label{yaa11p3e5}
\sum\limits_{n=1}^{\infty} na_n(z-z_0)^{n-1},
\end{equation}

\begin{equation}\label{yaa11p3e6}
\sum\limits_{n=0}^{\infty}\frac{a_n}{n+1}(z-z_0)^{n+1}
\end{equation}
имеют тот же радиус сходимости, что и ряд \eqref{yaa11p3e1}.
\end{thm}

\begin{proof}
Заметим, что ряд \eqref{yaa11p3e5} сходится тогда и только тогда, когда сходится ряд
$$
\sum\limits_{n=1}^{\infty}na_n(z-z_0)^n.
$$
Так как $\lim\limits_{n\to\infty} \sqrt[n]{n} = 1$, то последовательности $\{\sqrt[n]{|a_n|}\}$ и $\{\sqrt[n]{n|a_n|}\}$ имеют одни и те же частные пределы. Следовательно,
$$
\overline{\lim\limits_{n\to\infty}}\sqrt[n]{n|a_n|}=\overline{\lim\limits_{n\to\infty}}\sqrt[n]{|a_n|},
$$
что, в силу формулы Коши"--~Адамара, и доказывает равенство радиусов сходимости степенных рядов~\eqref{yaa11p3e1} и \eqref{yaa11p3e5}.

Аналогично ряд~\eqref{yaa11p3e6} сходится или расходится одновременно с рядом
$$
\sum\limits_{n=0}^{\infty}\frac{a_n}{n+1}(z-z_0)^n.
$$
А так как 
$$
\overline{\lim\limits_{n\to\infty}}\sqrt[n]{\frac{|a_n|}{n+1}}=\overline{\lim\limits_{n\to\infty}}\sqrt[n]{|a_n|},
$$
то ряд~\eqref{yaa11p3e6} имеет тот же радиус сходимости, что и ряд \eqref{yaa11p3e1}. Теорема доказана.

Отметим, что ряд~\eqref{yaa11p3e5} получается из ряда~\eqref{yaa11p3e1} почленным дифференцированием, а ряд~\eqref{yaa11p3e6} "--- почленным интегрированием по отрезку~$[z_0;z]$, соединяющему точки $z_0$ и $z$. Следовательно, теорему \eqref{yaa11p3thm3} можно сформулировать так:
\textit{Исходный ряд и ряды полученные из него почленным дифференцированием и почленным интегрированием, имеют один и тот же радиус сходимости.}
\end{proof}

\begin{thm}[Вторая теорема Абеля] \label{yaa11p3thm4}
Пусть $R$ "--- радиус сходимости степенного ряда \eqref{yaa11p3e1}. Тогда, если $R>0$, то ряд \eqref{yaa11p3e1} сходится равномерно на любом замкнутом круге $|z_0-z|\le r$, у которого $r<R$.
\end{thm}

\begin{cons}
Сумма степенного ряда непрерывна в круге сходимости.
\end{cons}

Рассмотрим теперь случай, когда степенной ряд сходится в некоторой точке, лежащей на границе его круга сходимости. Рассмотрим ряды вида

\begin{equation}\label{yaa11p32e2}
\sum\limits_{n=0}^{\infty} a_nz^n.
\end{equation}

\begin{thm}
Если $R$ "--- радиус сходимости ряда~\eqref{yaa11p32e2}, и ряд сходится при $z=R$, то он сходится равномерно на отрезке $[0;R]$ действительной оси.
\end{thm}

\begin{proof}
Если $z\in[0;R]$, то
$$
a_nx^n=a_nR^n\left(\frac{x}{R}\right)^n.
$$
Ряд $\sum\limits_{n=0}^{\infty} a_nR^n$ сходится равномерно относительно $x$, а последовательность
$$
b_n=\left(\frac{x}{R}\right)^n,\quad n\in \bbN,
$$
при любом $x$ монотонна и на $[0;R]$ ограничена:
$$
0\le \left(\frac{x}{R}\right)^n\le 1
$$
Поэтому, согласно принципу Абеля, ряд~\eqref{yaa11p32e2} сходится равномерно на $[0;R]$. 

Теорема доказана.
\end{proof}

Эта теорема называется \textit{Второй теоремой Абеля}.


\section{Бесконечная дифференцируемость. Ряд Тейлора. Аналитические функции}
\begin{thm}[Бесконечная дифферецнируемость суммы степенного ряда] \label{ch13.4thm0}
Пусть $f(x) = \sum\limits_{n=0}^{\infty} a_n (x - x_0)^n$ — сумма степенного ряда 
с радиусом сходимости $R \in (0; +\infty]$. Тогда функция $f$ бесконечно 
дифференцируема (т.е. имеет производные всех порядков $k = 1, 2, \ldots$) 
на интервале $(x_0 - R; x_0 + R)$; при этом при всех $k = 1, 2, \ldots$
\[
f^{(k)}(x) = \sum_{n=k}^{\infty} a_n \cdot n(n-1) \ldots (n - k + 1)(x - x_0)^{n-k},
\]
\[
x \in (x_0 - R;\, x_0 + R),
\]
и каждый такой ряд ($k$ раз формально продифференцированный исходный ряд) 
имеет тот же радиус сходимости $R$.
\end{thm}

\begin{proof}
Теорему достаточно доказать для $k = 1$, далее применять её последовательно. 
Формально продифференцированный ряд $\sum\limits_{n=1}^{\infty} n a_n (x - x_0)^{n-1}$ 
имеет тот же радиус сходимости, что и исходный, по лемме~17.2. Так как при всех 
$r \in (0; R)$ продифференцированный ряд равномерно сходится на отрезке 
$[x_0 - r;\, x_0 + r]$, а сам ряд сходится, например, в точке $x = x_0$, то по 
теореме о почленном дифференцировании функционального ряда равенство 
$f'(x) = \sum\limits_{n=1}^{\infty} n a_n (x - x_0)^{n-1}$ имеет место на отрезке 
$[x_0 - r;\, x_0 + r]$, а в силу произвольности $r \in (0; R)$ — и на всём 
интервале $(x_0 - R;\, x_0 + R)$. 
\end{proof}

\begin{defn}
Функция действительной переменной $f$ называется \emph{аналитической} в точке $x_0$, 
если существует $\delta > 0$ такое, что при всех $x \in (x_0 - \delta;\, x_0 + \delta)$ 
выполняется равенство
\[
f(x) = \sum_{n=0}^{\infty} a_n (x - x_0)^n
\]
(радиус сходимости последнего ряда $R \geqslant \delta$).
\end{defn}

\begin{thm}[Ряд Тейлора] \label{ch13.4thm1}
Если функция $f$ — аналитическая в точке $x_0$, то она имеет в точке $x_0$ 
производные всех порядков, причём коэффициенты степенного ряда, являющегося 
разложением функции $f$, вычисляются по формулам 
$a_k = \dfrac{f^{(k)}(x_0)}{k!}$, $k = 1, 2, \ldots$, т.е. разложение аналитической 
функции в степенной ряд в окрестности точки единственно:
\[
f(x) = \sum_{n=0}^{\infty} \frac{f^{(n)}(x_0)}{n!} (x - x_0)^n, 
\quad x_0 - \delta < x < x_0 + \delta
\]
(степенной ряд является рядом Тейлора функции $f$).
\end{thm}

\begin{proof}
Бесконечная дифференцируемость аналитической функции на всём интервале 
$(x_0 - \delta;\, x_0 + \delta)$ следует из теоремы~17.6; при этом $f(x_0) = a_0$. 
Так как при $k = 1, 2, \ldots$
\[
f^{(k)}(x) = \sum_{n=k}^{\infty} a_n \cdot n(n-1) \ldots (n - k + 1)(x - x_0)^{n-k},
\]
\[
x_0 - \delta < x < x_0 + \delta,
\]
то $f^{(k)}(x_0) = a_k \cdot k!$ (все слагаемые при $n > k$ обращаются в нуль), 
и $a_k = \dfrac{f^{(k)}(x_0)}{k!}$, $k = 1, 2, \ldots$ 
\end{proof}

\medskip
Обратное утверждение неверно: бесконечно дифференцируемая в окрестности точки 
функция не обязана быть аналитической в этой точке.

\begin{notion}
Рассмотрим функцию
\[
f(x) = 
\begin{cases}
e^{-1/x^2}, & \text{если } x \neq 0; \\
0, & \text{если } x = 0.
\end{cases}
\]
Для этой функции $f^{(k)}(0) = 0$ при всех $k = 0, 1, 2, \ldots$; график этой функции изображён на рис. 
Ряд Тейлора этой функции в точке $x = 0$ — 
нулевой. Сумма этого ряда равна $0$ при всех $x$, и ряд Тейлора функции $f$ 
не сходится к $f(x)$ ни в одной окрестности точки $0$ (т.е. эта функция 
бесконечно дифференцируема в любой точке, но не является аналитической в 
точке $x = 0$).
\end{notion}

\begin{thm}
Пусть функция $f(x)$ определена и бесконечно дифференцируема в некоторой $\delta$-окрестности точки $x_0$. Тогда, если
$$
\exists M\cquad |f^{(n)}(x)| \le M \quad \forall x \in O_\delta(x_0), \quad \forall n,
$$ 
то
$$
f(x) = \sum\limits_{n = 0}^{\infty} \frac{f^{(n)}(x_0)}{n!}(x - x_0)^n \quad \forall x \in O_\delta(x_0).
$$
\end{thm}
\begin{proof}
Для функции $f(x)$ напишем формулу Тейлора $(\ref{ch13.4eq1})$ с остаточным членом в форме Лагранжа:
$$
r_n(x) = \frac{f^{(n+1)}(\xi)}{(n + 1)!} (x - x_0)^{n + 1},
$$
где $\xi$ лежит между $x$ и $x_0$. Тогда $|f^{(n+1)}(\xi)| \le M$, и поэтому
$$
|r_n(x)| \le M \cdot \frac{\delta^{n + 1}}{(n + 1)!}
$$
для любого $n$ и любого $x \in O_\delta(x_0)$. А это и доказывает теорему, так как
\begin{equation*}
\lim\limits_{n \to \infty} \frac{\delta^{n + 1}}{(n + 1)!} = 0. \qedhere
\end{equation*}
\end{proof}
